
\documentstyle[12pt]{article}

\begin{document}

\begin{tabbing}
$2.0 \pm 0.6 \times 10^9$ gm/year \hspace*{0.5in}  \= Kane \& Gardner 1993
					 (``meteoric debris'') \= \kill
$1.5 \times 10^{11}$ gm/year \> Ceplecha 1996 \\
$4.0 \pm 2.0 \times 10^{10}$ gm/year \> Love \& Brownlee 1993 
					     (small particle flux only) \\
$2.0 \pm 0.6 \times 10^9$ gm/year \> Kane \& Gardner 1993
					 (``meteoric debris'') \\
$1.7 \times 10^{11}$ gm/year \> Ceplecha 1992 \\
$1.6 \times 10^9$ gm/year  \>        d'Almeida et al. 1991 \\
$7.8 \times 10^{10}$ gm/year  \>        Wasson \& Kyte 1987 \\
$1.6 \times 10^{10}$ gm/year  \>        Hughes 1978 \\
$2.0 \times 10^{10}$ gm/year  \>        Dohnanyi 1972 \\
\end{tabbing}

\noindent
Assume an infall rate of $1.5 \times 10^{8}$ kg/year (Ceplecha 1996).
\begin{tabbing}
$M_e$ \= = \= $1.48 \times 10^{14}$ m$^2$ \ \= \ \ [Area land surface] 
\= \ \ (ratio land/total = 0.290)\kill
$M_e$ \> = \> $5.973 \times 10^{24}$ kg \> [Mass of the earth] \\
$R_e$ \> = \> $6.371 \times 10^6$ kg    \> [Radius of the earth] \\
$A_e$ \> = \> $5.10 \times 10^{14}$ m$^2$ \> [Area total] \\
      \> = \> $1.48 \times 10^{14}$ m$^2$ \> [Area land surface] \> (ratio land/total = 0.290)\\
      \> = \> $3.62 \times 10^{14}$ m$^2$ \> [Area ocean surface] \> (ratio ocean/total = 0.710) \\
%      \> (ratio land/total = 0.290) \\
%      \> (ratio ocean/total = 0.710) \\
%      \> (ratio land/ocean = 0.409) \\
\end{tabbing}

\bigskip
\noindent
($1.5 \times 10^8$ kg/year) $\times$ ($4.5 \times 10^9$ year) = 
$6.75 \times 10^{17}$ kg

\bigskip
\noindent
${{6.75 \times 10^{17} \mbox{\tiny \  kg}} 
\over {5.973 \times 10^{24} \mbox{\tiny \  kg}}}$
= $1.130 \times 10^{-7}$

\bigskip
\noindent
${{6.75 \times 10^{20} \mbox{\tiny \ g}} \over {2.0 \mbox{\tiny \  g/cm}^3}}$
= $3.375 \times 10^{20}$ cm$^3$

\bigskip
\noindent
$A_e = 5.10 \times 10^{14} \mbox{\ m}^2 = 5.10 \times 10^{18} \mbox{\ cm}^2$



\end{document}